
\subsubsection{3.1 Experimental Design}

The core experiment holds constant all factors except task granularity: same model, same GRPO hyperparameters, same training duration. The function-level condition uses binary execution reward on competitive programming problems; the repository-level condition uses partial credit reward on software engineering tasks. Advantage variance is measured at every training step.

This design isolates the effect of reward sparsity on GRPO advantage variance. The function-level condition is expected to produce near-zero positive reward rates; the repository-level condition provides non-zero reward even for incomplete solutions through file-path overlap scoring.

\subsubsection{3.2 Model}

Both conditions use \textbf{CodeLlama-7b-Instruct-hf} [Rozière et al., 2023] (HuggingFace Hub: \texttt{codellama/CodeLlama-7b-Instruct-hf}) with no additional fine-tuning. The instruction-tuned variant is used to ensure the model can follow task formatting requirements. This is a 7B-class model with known near-zero solve rates on competitive programming problems at APPS difficulty.

Note: CodeLlama-7b-Instruct was used for the H-M1 advantage variance experiment. The intended model for the full curriculum GRPO experiment (H-E1) is DeepSeek-Coder-7B-base; the behavioral results from that model were not obtained.

\subsubsection{3.3 Datasets}

\textbf{Function-level condition:}
\begin{itemize}
\item \textbf{APPS} \citep{Hendrycksetal2021}: \texttt{codeparrot/apps}, train split, 5,000 problems with difficulty tiers 0–4. Problems require writing self-contained Python programs with stdin/stdout test cases.
\item \textbf{CodeContests} \citep{Lietal2022}: \texttt{deepmind/code_contests}, train split, 13,328 problems. Codeforces problems with public test cases.
\end{itemize}

\textbf{Repository-level condition:}
\begin{itemize}
\item \textbf{SWE-bench Verified} \citep{Jimenezetal2024}: \texttt{princeton-nlp/SWE-bench_Verified}, test split, 500 GitHub issues. Each task requires generating a unified diff patch resolving a real software engineering issue.
\end{itemize}

APPS and CodeContests use incompatible test case schemas (\texttt{input_output} string vs. \texttt{public_tests} dict). The implementation unifies these to a common 4-column format.

\subsubsection{3.4 GRPO Configuration}

Both conditions use identical hyperparameters:

\begin{table}[htbp]
\centering
\begin{tabular}{ll}
\toprule
Hyperparameter & Value \\
\midrule
Group size G & 8 \\
Batch size B & 4 prompts/step \\
Training steps & 120 \\
Temperature & 0.8 \\
Max new tokens & 512 \\
Framework & TRL 1.3.0 GRPOTrainer \\
Hardware & NVIDIA H100 NVL \\
\bottomrule
\end{tabular}
\end{table}

\subsubsection{3.5 Reward Functions}

\textbf{Function-level:} Binary execution reward. Each of G=8 completions is executed against the problem's test cases via subprocess (3-second timeout per test case). Reward = 1.0 if all test cases pass, 0.0 otherwise.

\textbf{Repository-level:} Partial credit based on file-path overlap. The generated unified diff is parsed for modified file paths. Reward = fraction of gold patch file paths appearing in the generated patch.

\subsubsection{3.6 Advantage Variance Measurement}

At each training step t, the per-step GRPO advantage variance is recorded:

$$\text{adv\_var}(t) = \text{Var}\left(\left\{ A_i^{(j)} : i \in [G], j \in [B] \right\}\right)$$

where A_i^(j) = (r_i^(j) − mean(r^(j))) / std(r^(j)). For degenerate groups where std(r^(j)) = 0, advantages are set to 0.0. This is logged by the \texttt{RewardDensityCallback} at every step, aggregating across all B×G = 32 advantage values per step.

\subsubsection{3.7 Statistical Analysis}

Per-step advantage variance between the two conditions is compared using Welch's two-sample t-test on log-transformed variance values. Log transformation is applied because advantage variance is right-skewed and bounded below at zero. Welch's t-test is used over Student's t-test because the two conditions are expected to have substantially different variances.

\begin{itemize}
\item Null hypothesis H₀: Mean advantage variance is equal between conditions.
\item Alternative hypothesis H₁: Mean advantage variance differs (two-tailed).
\end{itemize}

Effect size is reported as Cohen's d from log-transformed means and pooled standard deviation.

\subsubsection{3.8 Curriculum GRPO Infrastructure}

In parallel with the advantage variance experiment, training infrastructure for a 4-condition curriculum GRPO experiment is implemented and smoke-tested:

\begin{itemize}
\item \textbf{CurriculumDataset}: Wraps APPS+CodeContests with a \texttt{set_step()} interface that changes difficulty tier sampling at configurable step thresholds.
\item \textbf{CurriculumCallback}: Triggers difficulty tier transitions at pre-specified steps (easy tiers 0–2 for steps 0–2500, hard tiers 3–4 for steps 2501–5000 in the curriculum condition).
\item \textbf{RewardDensityCallback}: Logs per-step reward density and advantage variance to CSV files.
\end{itemize}

A 10-step smoke test verifies that all four conditions (curriculum, uniform, easy_only, hard_only) run without errors. Full 5000-step training was not executed.
